% Part: computability
% Chapter: recursive-functions
% Section: notations-pr-functions

\documentclass[../../../include/open-logic-section]{subfiles}

\begin{document}

\olfileid[hi]{cmp}{rec}{not}
\olsection{आदिम पुनरावर्तन की संकेत-पद्धतियाँ}

आदिम पुनरावर्ती फलनों का यथार्थ आगमनात्मक वर्णन होने का एक लाभ यह है कि
हम उनका व्यवस्थित ढंग से वर्णन कर सकते हैं। मसलन, हम प्रत्येक ऐसे फलन को
इस प्रकार एक ``संकेत'' दे सकते हैं। शून्य, उत्तराधिकारी और प्रक्षेपणों के
लिए क्रमशः $\Zero$, $\Succ$ और $\Proj{n}{i}$ चिह्नों का उपयोग करें। अब
मान लें कि $h$ को किसी $k$-स्थानी फलन~$f$ और $n$-स्थानी फलनों $g_0$,
\dots,~$g_{k-1}$ से संयोजन द्वारा परिभाषित किया गया है, और हमने बाद वाले
फलनों को संकेत $F$, $G_0$, \dots,~$G_{k-1}$ दिए हैं। तब एक नए चिह्न
$\fn{Comp}_{k,n}$ का उपयोग करके हम फलन $h$ को
$\fn{Comp}_{k,n}[F,G_0,\dots,G_{k-1}]$ से निरूपित कर सकते हैं।

आदिम पुनरावर्तन द्वारा परिभाषित फलनों के लिए हम अनुरूप संकेतों का उपयोग
कर सकते हैं। मान लें कि $(k+1)$-स्थानी फलन~$h$ को $k$-स्थानी फलन~$f$
और $(k+2)$-स्थानी फलन~$g$ से आदिम पुनरावर्तन द्वारा परिभाषित किया गया
है, और $f$ तथा~$g$ को क्रमशः $F$ और~$G$ संकेत दिए गए हैं। तब~$h$ को
दिया गया संकेत $\fn{Rec}_k[F,G]$ है।

याद करें कि जोड़ फलन को आदिम पुनरावर्तन द्वारा इस प्रकार परिभाषित किया गया है:
\begin{align*}
  \Add(x_0, 0) & = \Proj{1}{0}(x_0) = x_0\\
  \Add(x_0, y+1) & = \Succ(\Proj{3}{2}(x_0, y, \Add(x_0, y))) = \Add(x_0, y) +1
\end{align*}
यहाँ~$f$ की भूमिका $\Proj{1}{0}$ और~$g$ की भूमिका
$\Succ(\Proj{3}{2}(x_0, y, z))$ निभाता है। बाद वाले को संकेत
$\fn{Comp}_{1,3}[\Succ,\Proj{3}{2}]$ दिया जाता है, क्योंकि वह
$1$-स्थानी फलन~$\Succ$ और $3$-स्थानी फलन~$\Proj{3}{2}$ से संयोजन द्वारा
फलन परिभाषित करने का परिणाम है। इस व्यवस्था में हम जोड़ फलन को इस प्रकार
निरूपित कर सकते हैं:
\[
\fn{Rec}_1[\Proj{1}{0},\fn{Comp}_{1,3}[\Succ,\Proj{3}{2}]].
\]
ये संकेत कभी-कभी उपयोगी सिद्ध होते हैं, मसलन आदिम पुनरावर्ती फलनों को
परिगणित करते समय।

\begin{prob}
  ~$\Mult$ के लिए पूर्ण आदिम पुनरावर्ती संकेत दें।
\end{prob}

\end{document}
